\section{What the tracker must measure: the angle}\label{sec:requirement}
\sectionstate{partial}{new section (Dylan 2026-11-10: ``a diagram and section for arguing why angular resolution is most
important \ldots{} a quick toy simulation with visualization''). Toy and figure: \texttt{papers/paper1\_v2/figs/make\_pointing.py},
numbers in \texttt{figs/out/pointing.json}; scattering from \texttt{figs/out/material\_budget.json}. 2026-11-11: opening
angle and pointing separated, mass relation added, the origin-constraint paragraph redone with scattering.}
\intent{Before the detector is described: a chamber sees a short piece of track far from the target, so the angle, not
the position, decides both how well the opening angle of a pair is known (the mass) and how well a track points to its
origin (the vertex). Two requirements, one quantity; scattering sets the scale for both.}

A chamber does not see where a particle started. It sees a straight segment about 30\,mm long, 20\,cm away
(\crefp{fig:pointing}{a}), and everything the experiment wants to know about the track at the target has to be
extrapolated from that segment. Two measured quantities enter: the position of the segment, known to $\sigma_x$, and
its slope, known to $\sigma_\theta$. Carried back over the distance $d$ from the target to the plane where the position
is measured, the slope error moves the track by $d\,\sigma_\theta$ and the position error by $\sigma_x$, so the track
misses its origin by
\begin{equation}
  \sigma_b=\sqrt{\sigma_x^2+(d\,\sigma_\theta)^2}\,.\label{eq:pointing}
\end{equation}
With $d=233$\,mm from the target to the readout plane, an angle error of 1\textdegree{} is 4.1\,mm at the target.
A position error of 0.5\,mm is 1.5\,\% of the variance beside it, and the 0.12\,mm that the chambers reach
(\cref{sec:position}) is 0.1\,\%: the position would have to be as poor as 4\,mm to matter as much as one degree. The
lever arm is eight times the length of the segment, and that ratio is why the angle dominates. \Cref{fig:pointing}
shows it to scale, and a toy simulation that throws tracks from the target over the acceptance of a chamber, smears the
measured position and slope and extrapolates them back reproduces \cref{eq:pointing} (\crefp{fig:pointing}{b}; the toy
lies slightly above the line because inclined tracks have a longer lever arm).

The angle serves two purposes, and they are not the same requirement. The first is the opening angle $\theta$ of
the pair, which gives its mass. The two leptons share the energy $E$ of the transition, and with
$y=(E_+-E_-)/E$ and the electron mass neglected
\begin{equation}
  m^2=E^2\,(1-y^2)\,\sin^2(\theta/2)\,.\label{eq:mass}
\end{equation}
An X17 of 16.8\,MeV appears as an edge at $\theta=109$\textdegree{}, where one degree of opening angle is 0.10\,MeV
of mass. Each direction is measured inside its chamber, so the error on the opening angle is that of the two slopes.
The lever arm $d$ does not enter, and neither does the position. The second purpose is the pointing of
\cref{eq:pointing}. It tells whether the two tracks come from one point and whether that point lies in the target,
which rejects background but does not measure the mass. Both improve with $\sigma_\theta$; only the pointing depends
on the distance to the target.

What sets the useful scale for $\sigma_\theta$ is scattering before the measurement. In the 20\,cm of air between
the target and the window an 8\,MeV electron is deflected by 1.7\textdegree{} rms in each projection, and the window and
the cathode of the chamber add to it: 2.9\textdegree{} in all at 8\,MeV, 5.5\textdegree{} at 4\,MeV and 1.5\textdegree{} at
16\,MeV (\cref{tab:material}). That scattering is itself an angle error, of 4.5 to 17\,mm at the target (arrows in
\crefp{fig:pointing}{b}). It enters the opening angle in full: for a pair that shares the energy equally, with one
lepton in each of two chambers, a chamber resolution of 1\textdegree{} alone would give the opening angle to
1.3\textdegree{} (0.14\,MeV), and with the scattering it is 3.6\textdegree{} (0.38\,MeV).

A direction can also be taken from the position in the chamber and an assumed origin, and scattering affects that
direction less. A deflection in the foils just in front of the gas hardly moves the position, and one in the air moves
it by only part of the lever arm, so this direction is scattered by 1.2\textdegree{} at 8\,MeV where the local angle
is scattered by 2.9\textdegree. Its error is otherwise the uncertainty of the origin divided by $d$. Against a local
angle of 1\textdegree{} it would be better only for an origin known to 4\,mm. Against the local angle with its
scattering the limit is 21, 11 and 7\,mm at 4, 8 and 16\,MeV. The $^3$He volume is 20\,mm across and 60\,mm long:
about 5\,mm rms across the beam, inside the limit at every energy, and about 15\,mm along it, inside the limit only
for the softest leptons. The best estimate of the opening angle therefore combines the two, with the local angle
carrying the most weight for the hard leptons and along the beam. The position resolution is 0.03\textdegree{} over
this lever arm and limits neither.
\askDylan{New 2026-11-11 (your question on pointing against opening angle). The $^3$He bore, 20\,mm diameter and
40\,mm cylinder plus hemispherical caps, is from the Geant4 geometry (\texttt{MX17\_Full\_Geant/include/SimConfig.hh},
``per STEP''); the rms values assume captures fill it uniformly, which they do not (they sit towards the nose). Numbers:
\texttt{figs/out/pointing.json}, key \texttt{two\_estimators}.}

A tracker for these leptons therefore needs an angular resolution of about a degree, below the scattering at every
energy, a position resolution no better than a millimetre, and as little material as possible in front of its drift
volume. \Cref{sec:tracking} gives what the chambers deliver against this.

\begin{figure}[tbp]
  \centering
  \figorbox[width=0.9\linewidth]{pointing}{to-scale side view and toy pointing resolution}
  \caption{Why the angle. (a) One chamber to scale, seen from the side: the measured segment, 28\,mm long, lies
  233\,mm from the target. Inset, at the target and to scale: toy back-extrapolations for an angle error of
  1\textdegree{} with an exact position (blue), and for a position error of 0.5\,mm with an exact angle (orange).
  (b) Pointing resolution at the target ($s_{68}$ of the miss distance in one projection) against the angular
  resolution, for four position resolutions. Points: toy simulation over the acceptance of a chamber. Lines:
  \cref{eq:pointing} at normal incidence. Square: the chambers of this paper (\cref{sec:tracking}). Arrows: the
  pointing error from multiple scattering in the material in front of the drift gas (\cref{tab:material}) for electrons
  of 4, 8 and 16\,MeV.}
  \label{fig:pointing}
\end{figure}
